\section{Method}\label{sec:method}
\Cref{fig:pipeline} (\cref{app:data}) gives an overview. Every method sees the same cleaned, deduplicated and split data. The language
models read it through one of four text encodings, while XGBoost reads the table directly.


\paragraph{Encodings.} All encodings go through the features in the same fixed order, write missing values as
\texttt{NA} and join the parts with spaces (\cref{tab:examples}).
\textbf{E1 (PPFLE)} replaces each \texttt{column\$value} with its \res{cfg.encoders.hash.algorithm} hash, cut to
\res{cfg.encoders.hash.length} hexadecimal characters, as in SecurityBERT.
\textbf{E2 (binned PPFLE)} first maps each numeric value to one of \res{cfg.encoders.binning.n_bins} quantile bins,
fitted on the training split, and then hashes \texttt{column\$binK}. Categorical columns such as ports, protocols and
flags are hashed as in E1. \textbf{E3} writes readable \texttt{key=value} pairs with short column names. \textbf{E4}
fills a fixed template of English sentences and replaces protocol codes with their names.

\begin{table*}[!tbp]
  \centering
  \caption{The same NetFlow flow in the four encodings (beginning of each text only).}
  \label{tab:examples}
  % One flow in every encoding: generated by src/eval/analysis.py from results/encoding_stats.json, do not edit.
\begin{tabular}{@{}l p{0.78\linewidth}@{}}
\toprule
Encoding & Text \\
\midrule
E1 PPFLE & \texttt{\small e61bda9819b7fb55 4fbfe692dd218a7c 49f8efc3a07a7d29 46928dd8ca3a3f99 99a76dcb3e29cf31 da597cc9203f6df4 3efb85bbd25efb1a ece41bff11d60b21 3240b9e49a543820 fa2e50ac5f639e3f babb86cc04e6de22 c9436fb79e324b5e 0fce18341c7620a7 \dots} \\
E2 PPFLE-bins & \texttt{\small e61bda9819b7fb55 4fbfe692dd218a7c 49f8efc3a07a7d29 7fc09090edd30595 02e02bf5a981da22 37d4af979005d2cb 29fde00e2ab18ffe ece41bff11d60b21 3240b9e49a543820 fa2e50ac5f639e3f 34da606c4a74c86b 5b3a60bc305317a0 1126843ce2916164 \dots} \\
E3 key=value & \texttt{\small dst port=31992 proto=6 app proto=11.0 in bytes=4704 in pkts=28 out bytes=2976 out pkts=28 tcp flags=27 client flags=27 server flags=27 duration=335 in duration=334 out duration=334 min ttl=31 max ttl=32 longest pkt=984 shortest \dots} \\
E4 sentence & A TCP flow with application protocol NFS to destination port 31992 lasted 335 milliseconds. The client sent 28 packets (4704 bytes) over 334 milliseconds, and the server replied with 28 packets (2976 bytes) over 334 milliseconds. \dots \\
\bottomrule
\end{tabular}

\end{table*}

\paragraph{Models.} \textbf{XGBoost}, an ensemble of gradient-boosted decision trees, reads the cleaned table
directly. Its depth, learning rate and class weighting are chosen on the validation split, with early stopping.

Both language models are BERT-style Transformer encoders~\cite{devlin2019bert}. They read the token sequence of a
flow, and a classification layer on the representation of the first token predicts the class. The \textbf{small
BERT} starts from random weights and only ever sees flow texts. It follows the size stated for SecurityBERT: \res{cfg.bert_small.model.num_hidden_layers}
layers, hidden size \res{cfg.bert_small.model.hidden_size}, \res{cfg.bert_small.model.num_attention_heads} attention
heads and feed-forward size \res{cfg.bert_small.model.intermediate_size}. This gives \res{params.bert_small.millions}\,M
parameters, not the \res{exp1.paper.params_millions}\,M the paper states. It uses a byte-level BPE tokenizer with
\res{cfg.bert_small.tokenizer.vocab_size} tokens, trained separately for each encoding on the training split. The
maximum input length covers the \res{cfg.bert_small.tokenizer.coverage_percentile}th percentile of training rows. The
model is trained for \res{cfg.bert_small.training.epochs} epochs, and we keep the epoch that does best on validation.

\textbf{DistilBERT}~\cite{sanh2019distilbert} is a compressed version of BERT. It has half the layers of BERT-base and
was trained by knowledge distillation, that is, to reproduce BERT's predictions on the same English text (books and
Wikipedia). It keeps most of BERT's accuracy on language-understanding benchmarks while being smaller and faster.
Unlike the small BERT, it therefore starts with knowledge of English words and numbers, which is what H4 tests. We
fine-tune the pretrained uncased model (\res{params.distilbert.millions}\,M parameters) for
\res{cfg.distilbert.training.epochs} epochs with its own WordPiece tokenizer, which is not retrained. Hashes break into many WordPieces, so E1
needs \res{encoding.stats.nf_unsw_nb15.encodings.e1_ppfle.wordpiece_tokens.p99} tokens for nearly all rows. It is
therefore cut at DistilBERT's limit of \res{cfg.distilbert.tokenizer.max_len_cap} tokens. Side runs with shorter
hashes, where nothing is cut, show no systematic cost of this truncation (\cref{tab:side}, \cref{app:more}).

\paragraph{Tuning and protocol.} The learning rate of each language model is picked from a grid of
\res{cfg.bert_small.tuning.lr_grid.count} values on the validation split of one pilot dataset, NF-UNSW-NB15, and
reused for the other NetFlow datasets. Edge-IIoTset is tuned the same way on its own validation split.
% CHECK cfg.<model>.tuning.pilot.dataset and cfg.<model>_edge.tuning.pilot.dataset
Every method sees the same samples, splits and columns, and the test split is used only once. Every model is run
with \res{seeds.distilbert} seeds and reported as mean $\pm$ standard deviation. DistilBERT is the most expensive
model, so its second and third seeds were trained on an A100 GPU of the Elja cluster, with the same settings and
learning rates as the first. Only the two truncation side runs use a single seed.
% CHECK seeds.xgboost = seeds.bert_small = seeds.distilbert

Macro-F1 is our main metric. We also report accuracy, to compare with the paper, and for cross-dataset runs the
attack recall and false-positive rate (FPR). For the hypotheses, we treat two numbers as different only if they
differ by more than \res{cfg.analysis.hypotheses.tie_margin} and by more than the larger of their seed standard
deviations. Otherwise we call them tied. We call an encoding ``good'' if it is within
\res{cfg.analysis.hypotheses.well_margin} of the best method.

